\section{Reproducibility}
\label{sec:reproducibility}

Every figure in this report is reproducible from the repository at a single
identified revision. The full protocol is documented in
\texttt{REPRODUCING.md}; the short version follows.

\paragraph{One command.}
\texttt{./scripts/run\_tournament.sh} runs five phases in order and refuses
to time anything before equivalence passes:

\begin{enumerate}
    \item \textbf{equivalence} --- cross-language canonical files, compared
    byte for byte (\texttt{equivalence\_\allowbreak report.json});
    \item \textbf{benchmarks} --- Go, Rust, Java (+JMH) and Python suites,
    each writing its CSVs under
    \texttt{benchmarks\_\allowbreak results/\allowbreak <language>/};
    \item \textbf{verify} --- schema check of every CSV
    (\texttt{scripts/\allowbreak verify\_\allowbreak format.py});
    \item \textbf{environment} --- host, runtimes and Git revision in
    \texttt{benchmarks\_\allowbreak results/\allowbreak environment.txt};
    \item \textbf{charts} --- \texttt{src/\allowbreak visualizations/\allowbreak main.py} renders \texttt{charts/}.
\end{enumerate}

There is no resume: a failed run restarts clean, so the committed snapshot
always comes from one continuous execution.

\paragraph{Chart-to-artifact map.}
\Cref{tab:chart-map} links every figure to its script (all chart scripts
live under \texttt{src/\allowbreak visualizations/\allowbreak modules/}), its raw results and its
configuration.

\begin{table}[htbp]
    \centering
    \small
    \caption{Where each figure comes from.}
    \label{tab:chart-map}
    \begin{tabularx}{\textwidth}{@{}llX@{}}
        \toprule
        Figure & Script & Raw results \\
        \midrule
        \cref{fig:radar} & \texttt{radar\_chart.py} & all of \texttt{benchmarks\_results/} \\
        \cref{fig:throughput} & \texttt{throughput\_charts.py} & \texttt{\*/throughput.csv}, \texttt{\*/benchmarks.csv} \\
        \cref{fig:latency} & \texttt{latency\_charts.py} & \texttt{\*/statistics.csv} \\
        \cref{fig:scalability-elapsed,fig:scalability-memory,fig:updates} & \texttt{scalability\_charts.py} & \texttt{\*/scalability.csv}, \texttt{\*/updates.csv} \\
        Sankey (interactive HTML) & \texttt{storage\_sankey.py} & \texttt{\*/disk\_usage.csv} \\
        \cref{tab:jmh} & \texttt{multilang\_ports/java/run\_jmh.sh} & \texttt{java/jmh\_results.json} \\
        \bottomrule
    \end{tabularx}
\end{table}

\paragraph{Revision.}
The Git commit that produced the committed snapshot is recorded inside
\texttt{benchmarks\_\allowbreak results\_\allowbreak environment.txt} together with the host and
runtime versions of \cref{tab:host}. Checking out that commit and re-running
the command above reproduces the pipeline; the raw numbers will differ in
their last digits (shared-host jitter, \cref{sec:limitations}), but the
ordering and the equivalence digests are stable.
